%======================================================================
\subsection{BW-SIM-06 phase 1 --- Built, 31 Oct: airflow with OpenFOAM, external flow}\label{sec:sim06-built}
%======================================================================

Phase 1 of \S\ref{sec:sim06-plan} is in (commit \code{SIM-06: airflow with OpenFOAM, phase 1}): the case writer and
runner in \file{src/Mod/BruceWorx/bruceworx\_cfd.py} (no FreeCAD needed to test it), the panels in
\file{bruceworx\_cfd\_gui.py}, and small hooks in the study frame (\file{bruceworx\_sim.py}, \file{bruceworx\_sim\_gui.py}).
Your answers O1--O9 were followed; nothing is installed by Vector Labs.

\paragraph{What you can do now.}
\begin{itemize}[nosep]
  \item \textbf{New Study $\to$ Airflow $\to$ External flow}: speed, angle of attack, reference area (blank = the
        frontal area of the bounding box), mesh preset (coarse, medium, fine), turbulence ($k$--$\omega$ SST or
        laminar), air temperature and altitude. The free stream is always $+X$ of the model; the angle of attack
        turns the model about Y, so drag is along $+X$ and lift along $+Z$.
  \item \textbf{Pre-check} finds \file{OpenFOAM-v2606.app} (\file{/Applications}, \file{\textasciitilde/Applications}, or the
        path in Preferences $\to$ Simulation $\to$ Airflow), warns on any other version, and when none is found fails
        with an \emph{Install OpenFOAM\ldots} button (opens the openfoam-app page; the text also gives
        \code{brew install gerlero/openfoam/openfoam@2606}). It also checks that every part is watertight, estimates
        the cell count and $y^+$, and warns when the Reynolds number points to the other turbulence choice.
  \item \textbf{Run} writes the case under \file{\textasciitilde/Library/Caches/Vector Labs/cfd/<document>/<study>/}
        (one STL per part, \code{blockMesh} box from the bounding box with editable upstream, downstream and sideways
        factors, \code{snappyHexMesh}, \code{simpleFoam}, forces and force coefficients per part, the inlet and outlet flow
        rate) and runs it in the background with a live residual plot, a live drag plot and a Stop button.
  \item \textbf{Results}: a forces table (drag, side, lift, $C_d$, $C_l$ per part and in total), the convergence verdict
        in the plan's three tests, \emph{Open in ParaView} when ParaView.app exists, and \emph{Report}, which adds an
        Airflow section (setup, convergence, forces) to the LaTeX report.
\end{itemize}

\paragraph{Benchmarks} (coarse preset, this Mac, M-series, one core, see tests \code{TestBruceWorxCfd}).
\begin{center}\small
\begin{tabular}{@{}llll@{}}
\toprule Case & Result & Reference & Time\\ \midrule
Laminar pipe, $D=10$\,mm, 0.5\,m/s & $\mathrm dp/\mathrm dx=2.905$\,Pa/m & 2.901 (plan 2.91), $+0.12$\,\% & 21\,s\\
Sphere, $Re=100$ (130\,k cells) & $C_d=1.110$ & 1.092 (plan 1.09), $+1.6$\,\% & 24\,s\\
\bottomrule
\end{tabular}
\end{center}
Neither is slow enough for a slow marker. A mixed three-part test (body and two arms, 5$^\circ$ incidence, $k$--$\omega$
SST with wall layers, 70\,k cells) converged in 172 iterations and 17\,s.

\paragraph{Decided while building.}
\begin{itemize}[nosep]
  \item \textbf{No MPI here.} \code{mpirun} hangs in the PMIx handshake of OpenFOAM-v2606's Open MPI 5 on this Mac
        (every \code{-parallel} start, with or without the sandbox, with the interface options tried). The run script
        starts \code{decomposePar} and \code{mpirun} on the performance cores, and if \code{simpleFoam} has written nothing
        after 25\,s it falls back to one core and the study remembers it for a day. The parallel path itself (its
        success branch with \code{reconstructPar}) has therefore not run yet.
  \item \textbf{Spaces in the folder.} OpenFOAM aborts on a case path with a space, and the cache folder is
        \file{Vector Labs}; the runner works through a space-free symlink in the temp folder.
  \item \textbf{Convergence.} The solver stops at residuals of $10^{-5}$, ten times under the limit, so the report's
        $10^{-4}$ holds with room. A run stopped early is judged on the forces of its last half (at least 50
        iterations) instead of 200, so the start-up transient does not count.
  \item \textbf{Results in 3D.} The headless path loads \code{foamToVTK} output into FreeCAD's FEM post pipeline
        (scaled to mm) with a cut plane. In the GUI build, creating any \code{Fem::FemPostPipeline} crashes
        (\code{SoFCColorGradient::setMarkerLabel}, known since 27 Oct), so the GUI converts to VTK and relies on ParaView
        until that is fixed. This build has no streamline filter either.
\end{itemize}
\paragraph{Open.} Streamlines and surface pressure inside Vector Labs wait on the pipeline crash. The preset cell counts
(about 0.3\,M, 1\,M, 3\,M for a drone-sized model) are estimates until measured on the SP-01 frame in phase 2.
